\scanpage{141}
% archive foliation: 140
\margin{$\deleted{\unsure{P_{\mathcal D}}}$}

Nous pouvons expliciter le fibré tangent le long \unsure{de|des} \deleted{fibres}\inserted{la section nulle} de la courbe elliptique universelle rigidifiée à la Torelli, comme étant le fibré formé des $t_E$, qui du point de vue fibré vectoriel réel s'identifie au \deleted{fibré} $\mathrm{Sl}(2,\mathbb Z)$-fibré sur $\mathcal D$ défini par l'opération \unsure{standart|standard} de $\mathrm{Sl}(2,\mathbb Z)$ sur $\mathbb R^2$. Le fibré principal associé \inserted{$P_{\mathcal D}$} s'obtient donc (pour sa seule structure topologique) en remplaçant dans cet énoncé $\mathbb R^2$ par $(\mathbb R^2)^*$.

Pour en trouver un revêtement universel, il suffit de choisir un revêtement universel de $(\mathbb R^2)^*\simeq\mathbb C^*$, p.ex. le revêtement par l'exponentielle complexe
\[
z\longmapsto\exp(2i\pi z):\quad
\mathbb C\longrightarrow\mathbb C^*.
\]
